\begin{table}[H]
\centering
\caption{Βασικά Στοιχεία \ENG{Dublin} \ENG{Core} \ENG{Metadata}}
\label{tab:dublin_core_elements}
\small
\setlength{\tabcolsep}{5pt}
\renewcommand{\arraystretch}{1.25}
\begin{chaptertablebox}
\begin{tabular}{|l|p{8cm}|}
\hline
\textbf{Στοιχείο} & \textbf{Περιγραφή} \\
\hline
\hline
\ENG{Title} & Το όνομα που δίνεται στον πόρο \\
\ENG{Creator} & Η οντότητα που είναι κυρίως υπεύθυνη για τη δημιουργία του πόρου \\
\ENG{Subject} & Το θέμα του περιεχομένου του πόρου \\
\ENG{Description} & Μια περιγραφή του περιεχομένου του πόρου \\
\ENG{Publisher} & Η οντότητα που είναι υπεύθυνη για τη διάθεση του πόρου \\
\ENG{Contributor} & Οντότητα που είναι υπεύθυνη για συνεισφορές στο περιεχόμενο \\
\ENG{Date} & Ημερομηνία που σχετίζεται με το κύκλο ζωής του πόρου \\
\ENG{Type} & Η φύση ή το είδος του περιεχομένου \\
\ENG{Format} & Το φυσικό ή ψηφιακό μορφότυπο του πόρου \\
\ENG{Identifier} & Μια μοναδική αναφορά στον πόρο (π.χ. \ENG{URI}, \ENG{ISBN}) \\
\ENG{Source} & Πόρος από τον οποίο προέρχεται ο παρών πόρος \\
\ENG{Language} & Η γλώσσα του περιεχομένου του πόρου \\
\ENG{Relation} & Αναφορά σε σχετικό πόρο \\
\ENG{Coverage} & Το χωρικό ή χρονικό πεδίο εφαρμογής του πόρου \\
\ENG{Rights} & Πληροφορίες για δικαιώματα που σχετίζονται με τον πόρο \\
\hline
\end{tabular}
\end{chaptertablebox}
\end{table}

